\documentclass[10pt,a4paper]{amsart}
\usepackage[margin=19mm]{geometry}
\usepackage{amsmath,amssymb,mathtools}
\usepackage[scr=rsfs,cal=euler]{mathalfa}
\usepackage{xfrac}
\usepackage[unicode,hidelinks,bookmarks=false]{hyperref}
\newcommand{\ie}{i.e.\ }
\newcommand{\cf}{cf.\ }
\newcommand{\resp}{respectively\ }
\newcommand{\Subsection}{\subsection}
\providecommand{\nobreakdash}{\nobreak\hbox{-}}
\setlength{\emergencystretch}{2em}
\multlinegap0pt
\usepackage{amssymb}
\usepackage{subfig}
\usepackage{tikz}
\usepackage{booktabs}
\usepackage{multirow}
\usepackage{bm}
\usepackage{float}
\usepackage{natbib}
\newtheorem{theorem}{Theorem}
\newtheorem{lemma}{Lemma}
\title{Martingale Characterizations of Non-Homogeneous Counting Processes and Their Fractional Variants}
\author{Kartik Tathe, Sayan Ghosh}
\date{Source: arXiv:2511.14328v2 (CC BY 4.0)}
\begin{document}
\maketitle

\noindent\emph{Source and licence.} Martingale Characterizations of Non-Homogeneous Counting Processes and Their Fractional Variants. Version arXiv:2511.14328v2 (CC BY 4.0), distributed by arXiv under the Creative Commons Attribution 4.0 International licence, http://creativecommons.org/licenses/by/4.0/, as recorded in the arXiv licence metadata for this exact version and shown on the abstract page https://arxiv.org/abs/2511.14328v2. Full licence notice: \texttt{licenses/arxiv-2511.14328-CC-BY-4.0.txt}.\par\smallskip

\noindent\textit{Editorial scope.} Counted unit: the article's theorem Watanabe\_characterization\_NTSTFPP (source0.tex lines 825--835). The article's complete original proof of it is delivered with it (source0.tex lines 836--890, one \texttt{proof} environment of 55 lines). No mathematics is written, repaired or re-derived: the statement and the proof are the article's own lines, in the article's own notation, and no retained line is substituted or re-flowed.  Retained uncounted as the source-local context the proof needs: lines 431--433; lines 730--732; lines 734--744.  Nothing was omitted from any retained span, so the package carries no omission record.  No retained line contains a citation, so the wrapper carries no bibliography and no bibliography entry.  No retained line contains a figure environment, an image-inclusion command or drawing code, so no image is delivered and the package declares no asset dependency.  Editorial formatting only, in addition: an AMS \texttt{amsart} preamble that restates the article's own declarations and macros used by the retained lines, namely the environments \texttt{theorem}, \texttt{lemma} on separate counters, together with the standard packages those lines need.  The retained environments print in this document as Theorem 1, Lemma 1 in order of appearance, whereas the article prints the same items with the numbering of its own full document.  The complete native source of the article as distributed in the arXiv e-print for this exact version is bundled unchanged as \texttt{sources/189352/source0.tex}.
\begin{theorem}\label{watanabe_npp}
Let $\{\mathcal{N}(t)\}_{t\ge 0}$ be a simple locally finite point process adapted to $\mathcal{F}_t$ where $\{\mathcal{F}_t=\sigma(\mathcal{N}(t))\}_{t\ge 0}$ is the natural filtration of $\{\mathcal{N}(t)\}_{t\ge 0}$. Then $\{\mathcal{N}(t)-\Lambda(t)\}_{t\ge 0}$ is a $\mathcal{F}_t$-martingale for some $\Lambda(t)>0$ if and only if $\{\mathcal{N}(t)\}_{t\ge 0}$ is a NPP with cumulative intensity function $\Lambda(t)$. 
\end{theorem}

\begin{lemma}\label{martingale_lemma}
    Let $X$ be a right-continuous martingale. If $T$ and $S$ are stopping times such that $P(T<\infty)=1$ and $\{X(t\wedge T),t\ge 0\}$ is uniformly integrable, then $\mathbb{E}(X(T)|\mathcal{F}_{S\wedge T})=X(S\wedge T)$.
\end{lemma}

\begin{theorem}\label{Watanabe_characterization_NFPP}
    Let $\{K(t)\}_{t\ge 0}$ be a simple locally finite point process. Then $\{K(t)\}_{t\ge 0}$ is a NTFPP if and only if there exist a rate function $\lambda:[0,\infty)\to[0,\infty)$ with $\Lambda(t)=\int_{0}^{t}\lambda(u)du<\infty$ and an inverse stable subordinator $\{Y_{\alpha}(t)\}_{t\ge 0}$ such that the process 
    \begin{equation}\label{martingale_form_nfpp}
    \{X(t)\}_{t\ge 0}=\{K(t)-\Lambda(Y_{\alpha}(t))\}_{t\ge 0}    
    \end{equation} 
    is a right-continuous martingale with respect to the induced filtration $\mathcal{F}_t=\sigma\left(K(s),s\le t\right)\vee\sigma\left(Y_{\alpha}(s), s\ge 0\right)$ and for any $T>0$,
    \begin{equation}\label{stopping_times_martingale_nfpp}
        \{X(\tau),\,\tau\,\, \text{is stopping time such that}\,\,Y_{\alpha}(\tau)\le T\}
    \end{equation}
    is uniformly integrable.
\end{theorem}

\begin{theorem}\label{Watanabe_characterization_NTSTFPP}
    Let $\{L(t)\}_{t\ge 0}$ be a simple locally finite point process. Then $\{L(t)\}_{t\ge 0}$ is NTSTFPP if and only if there exist a rate function $\lambda:[0,\infty)\to[0,\infty)$ with $\Lambda(t)=\int_{0}^{t}\lambda(u)du<\infty$, a tempered stable subordinator (TSS) $\{D_{\beta,\theta}(t)\}_{t\ge 0}$ and an inverse stable subordinator $\{Y_\alpha(t)\}_{t\ge 0}$ which are independent of each other such that the process 
    \begin{equation*}\label{martingale_form_NTSTFPP}
    \{R(t)\}_{t\ge 0}=\{L(t)-\Lambda(D_{\beta,\theta}(Y_{\alpha}(t)))\}_{t\ge 0}    
    \end{equation*} 
    is a right-continuous martingale with respect to the induced filtration $\mathcal{F}_t=\sigma\left(L(s),s\le t\right)\vee\sigma\left(D_{\beta,\theta}(Y_{\alpha}(s)), s\ge 0\right)$ and for any $T>0$,
    \begin{equation*}\label{stopping_times_martingale_NTSTFPP}
        \{R(\tau),\,\tau\,\, \text{is stopping time such that}\,\,D_{\beta,\theta}(Y_{\alpha}(\tau))\le T\}
    \end{equation*}
    is uniformly integrable.
\end{theorem}

\begin{proof}   
The proof is similar to that of Theorem \ref{Watanabe_characterization_NFPP} with $Y_{\alpha}(t)$ replaced by $D_{\beta,\theta}(Y_{\alpha}(t))$ and hence omitted.
% Let $\{L(t)\}_{t\ge 0}$ be a NTSTFPP so that $L(t)=\mathcal{N}\left(D_{\beta,\theta}(Y_{\alpha}(t))\right)$, where $\{D_{\beta,\theta}(t)\}_{t\ge 0}$ is a tempered stable subordinator and $\{Y_\alpha(t)\}_{t\ge 0}$ is an inverse stable subordinator which are independent of each other and also independent of the NPP $\{\mathcal{N}(t)\}_{t\ge 0}$ with rate function $\lambda(t)$.
%     For the NPP $\{\mathcal{N}(t)\}_{t\ge 0}$, the compensated process $\{\mathcal{N}(t)-\Lambda(t)\}_{t\ge 0}$
% is a right-continuous martingale with respect to the natural filtration of \(\mathcal{N}(t)\). Define
% \begin{equation*}
%     R(t):=L(t)-\Lambda\left(D_{\beta,\theta}(Y_{\alpha}(t))\right)=\mathcal{N}\left(D_{\beta,\theta}(Y_{\alpha}(t))\right)-\Lambda\left(D_{\beta,\theta}(Y_{\alpha}(t))\right).
% \end{equation*}
% First we show $\{R(t)\}_{t\ge 0}$ is a martingale with respect to $\mathcal{F}_t=\sigma\left(L(s),s\le t\right)\vee\sigma\left(D_{\beta,\theta}(s), s\ge 0\right)$.

% For \(0\le s<t\), using the tower property and conditioning on $D_\beta(t)$, we have
% \begin{equation*}
%     E\left[R(t)\mid\mathcal F_s\right]
% =E\left[\,E\left[\mathcal{N}\left(D_{\beta,\theta}(Y_{\alpha}(t))\right)-\Lambda\left(D_{\beta,\theta}(Y_{\alpha}(t))\right)\;\big|\;D_{\beta,\theta}(Y_{\alpha}(t))\vee\mathcal F_s\right]\;\Big|\;\mathcal F_s\right].
% \end{equation*}
% Given the path of $D_{\beta,\theta}(Y_{\alpha}(t))$, \(\mathcal{N}(D_{\beta,\theta}(Y_{\alpha}(t)))\) is a non-homogeneous Poisson process. Using the martingale property of the compensated non-homogeneous Poisson process, we get
% \[
% E\big[\mathcal{N}\big(D_{\beta,\theta}(Y_{\alpha}(t))\big)-\Lambda\big(D_{\beta,\theta}(Y_{\alpha}(t))\big)\;\big|\;D_{\beta,\theta}(Y_{\alpha}(t))\vee\mathcal F_s\big]
% = \mathcal{N}\big(D_{\beta,\theta}(Y_{\alpha}(s))\big)-\Lambda\big(D_{\beta,\theta}(Y_{\alpha}(s))\big).
% \]
% Therefore
% \[
% E\big[R(t)\mid\mathcal F_s\big]=\mathcal{N}\big(D_{\beta,\theta}(Y_{\alpha}(s))\big)-\Lambda\big(D_{\beta,\theta}(Y_{\alpha}(s))\big)=R(s),
% \]
% which proves that $\{R(t)\}_{t\ge 0}$ is a martingale. Right-continuity follows from right-continuity of $\mathcal{N}(t)$ and of $D_{\beta,\theta}(Y_{\alpha}(t))$.
   
%  Next we show that the family \(\{R(\tau):D_{\beta,\theta}(Y_{\alpha}(\tau))\le T\}\) is uniformly integrable. Fix \(T>0\) and let \(\tau\) be any stopping time with $D_{\beta,\theta}(Y_{\alpha}(\tau))\le T$. 
%  Note that
%  \begin{equation}\label{mean_ntstfpp_martingale}
%      \mathbb{E}\left[R(\tau)\right]=\mathbb{E}\left[\mathcal{N}(D_{\beta,\theta}(Y_{\alpha}(\tau)))-\Lambda(D_{\beta,\theta}(Y_{\alpha}(\tau)))\right]=\mathbb{E}\left[\Lambda(D_{\beta,\theta}(Y_{\alpha}(\tau)))\right]-\mathbb{E}\left[\Lambda(D_{\beta,\theta}(Y_{\alpha}(\tau)))\right]=0.
%  \end{equation}
%  Using the law of total variance, we have from \eqref{martingale_form_NTSTFPP}
% \begin{align*}
%     \mathbb{V}\left(R(\tau)\right)&=\mathbb{E}\left(\mathbb{V}\left[L(\tau)-\Lambda(D_{\beta,\theta}(Y_{\alpha}(\tau)))\mid D_{\beta,\theta}(Y_{\alpha}(\tau))\right]\right)
% +\mathbb{V}\left(\mathbb{E}[L(\tau)-\Lambda(D_{\beta,\theta}(Y_{\alpha}(\tau)))\mid D_{\beta,\theta}(Y_{\alpha}(\tau))]\right)\\
%     \implies\mathbb{V}\left(R(\tau)\right)&=\mathbb{E}\left(\mathbb{V}\left[\mathcal{N}(D_{\beta,\theta}(Y_{\alpha}(\tau)))\mid D_{\beta,\theta}(Y_{\alpha}(\tau))\right]\right)
% +\mathbb{V}\left(\mathbb{E}[\mathcal{N}(D_{\beta,\theta}(Y_{\alpha}(\tau)))\mid D_{\beta,\theta}(Y_{\alpha}(\tau))]\right)\\
%     \implies\mathbb{E}\left[R(\tau)^2\right]
% &=\mathbb{E}\left(\mathbb{V}\left[\mathcal{N}(D_{\beta,\theta}(Y_{\alpha}(\tau)))\mid D_{\beta,\theta}(Y_{\alpha}(\tau))\right]\right)
% +\mathbb{V}\left(\mathbb{E}[\mathcal{N}(D_{\beta,\theta}(Y_{\alpha}(\tau)))\mid D_{\beta,\theta}(Y_{\alpha}(\tau))]\right)+\left(\mathbb{E}\left[R(\tau)\right]\right)^2\\
% &=\mathbb{E}\left(\mathbb{V}\left[\mathcal{N}(D_{\beta,\theta}(Y_{\alpha}(\tau)))\mid D_{\beta,\theta}(Y_{\alpha}(\tau))\right]\right)
% +\mathbb{V}\left(\mathbb{E}[\mathcal{N}(D_{\beta,\theta}(Y_{\alpha}(\tau)))\mid D_{\beta,\theta}(Y_{\alpha}(\tau))]\right)\quad(\text{using}~\eqref{mean_ntstfpp_martingale})\\
% &=\mathbb{E}\left(\Lambda[D_{\beta,\theta}(Y_{\alpha}(\tau))]\right)
% +\mathbb{V}\left(\Lambda[D_{\beta,\theta}(Y_{\alpha}(\tau))]\right)\\
% &=\mathbb{E}\left(\Lambda[D_{\beta,\theta}(Y_{\alpha}(\tau))]\right)
% +\mathbb{E}\left(\left(\Lambda[D_{\beta,\theta}(Y_{\alpha}(\tau))]\right)^2\right)-\left(\mathbb{E}\left(\Lambda[D_{\beta,\theta}(Y_{\alpha}(\tau))]\right)\right)^2\\
% &\le \Lambda(T)+\left(\Lambda(T)\right)^2.
% \end{align*}
 
% Therefore, the family \eqref{stopping_times_martingale_NTSTFPP} is uniformly bounded in $L^2$, and hence uniformly integrable.
%    Conversely, it is enough to prove that $L(t)=\mathcal{N}\left(D_{\beta,\theta}(Y_{\alpha}(t))\right)$, where $\{\mathcal{N}(t)\}_{t\ge 0}$ is a NPP, independent of the subordinator $\{D_{\beta,\theta}(Y_{\alpha}(t))\}_{t\ge 0}$ with rate $\lambda(t)>0$. Let $Z(t)=\inf\{s:D_{\beta,\theta}(Y_{\alpha}(s))\ge t\}$ be the inverse of $\{D_{\beta,\theta}(Y_{\alpha}(t))\}_{t\ge 0}$. It can be observed that $\{Z(t),t\ge 0\}$ forms a family of stopping times. Therefore, by Lemma \ref{martingale_lemma}, 
%    \begin{equation*}
%        R\left(Z(t)\right)=L(Z(t))-\Lambda\left( D_{\beta,\theta}(Y_{\alpha}(Z(t)))\right)
%    \end{equation*} is still a martingale. As $D_{\beta,\theta}(Y_{\alpha}(.))$ is continuous, we have $D_{\beta,\theta}(Y_{\alpha}(Z(t)))=t$ which implies that $\{L(Z(t))-\Lambda(t)\}_{t\ge 0}$ is a martingale. Furthermore, as $Z(t)$ is an increasing process, $L(Z(t))$ is a simple point process. By Theorem \ref{watanabe_npp}, it follows that $L(Z(t))$ is a classical non-homogeneous Poisson process with rate $\lambda(t)>0$. If we define $\mathcal{N}(t)=L(Z(t))$, then the process $\{L(t)=\mathcal{N}(D_{\beta,\theta}(Y_{\alpha}(t)))\}_{t\ge 0}$ represents a NTSTFPP. This proves the theorem.
\end{proof}

\end{document}
